\section{VALIDAÇÃO DA MEDIÇÃO}

Antes de usar o sentimento por aspecto para dizer algo sobre a nota, é preciso saber o quanto essa medida merece confiança. Esta seção reúne cinco tipos de evidência, da mais direta à mais indireta: frases rotuladas por um pesquisador, concordância entre os dois modelos, uma regra fixada de antemão para escolher o modelo principal, concordância com as subnotas oficiais e uma checagem com um segundo método de tópicos e com um segundo anotador. A seção termina com uma classificação dos aspectos, que as Seções 5 a 7 respeitam.

\subsection{Evidência direta: frases rotuladas}

Um pesquisador rotulou, sem ver as classificações automáticas, 297 sentenças, indicando o sentimento (negativo, neutro ou positivo) e o aspecto principal. A amostra foi estratificada pelo aspecto e pela polaridade previstos pelo modelo A, com cerca de um terço de sentenças de cada polaridade por aspecto. Por isso a amostra é rica em sentenças difíceis: no corpus, 78,9\% das sentenças são positivas segundo A. Cada métrica é lida de duas formas, na amostra e ponderada pela população, em que cada sentença pesa o tamanho do seu estrato original (aspecto e polaridade previstos por A). Os pesos são os mesmos para A e B. O tamanho efetivo da amostra ponderada é de apenas 72 sentenças. Os intervalos de confiança são, em consequência, largos.

A Tabela~\ref{tab:rotuladas} apresenta as métricas de sentimento. Na amostra, o modelo B tem acurácia maior (63,3\% contra 58,9\%). Ponderando pela população, o modelo A é superior, com acurácia de 87,2\% contra 84,9\% e kappa ponderado linear de 0,66 contra 0,51.

\begin{table}[htb]
\centering
\caption{Sentimento: desempenho dos modelos A e B nas 297 sentenças rotuladas.}
\small
\begin{tabular}{l|r|r}
\hline
Medida & A & B\\
\hline
Acurácia na amostra & 58,9\% & 63,3\%\\
$F_1$ macro na amostra & 0,57 & 0,52\\
Kappa de Cohen na amostra & 0,39 & 0,38\\
Kappa ponderado linear na amostra & 0,48 & 0,50\\
Correlação de postos, rótulo humano $\times$ estrelas esperadas & 0,64 & 0,63\\
Acurácia ponderada pela população & 87,2\% & 84,9\%\\
Kappa ponderado pela população & 0,55 & 0,43\\
Kappa ponderado linear pela população & 0,66 & 0,51\\
\hline
\end{tabular}
\label{tab:rotuladas}
\end{table}

A discordância entre a amostra e a população tem explicação. A amostra foi construída a partir das previsões de A e concentra as sentenças em que A erra, sobretudo as que A chama de neutras. Nos dois modelos o erro mais comum é classificar como neutras sentenças que o pesquisador considera positivas: quando A prevê positivo, acerta em 94,9\% dos casos, mas quando prevê neutro acerta em apenas 26,1\%. O modelo B classifica mais sentenças como positivas (151 contra 98 das 297) e acerta menos nas previstas como neutras (16,9\%) e negativas (51,7\%, contra 61,9\% de A). Por idioma, o português, que responde por 66\% das avaliações, é o de menor acerto na amostra para o modelo A (51,8\%, contra 61,4\% de B); ponderada pela população, a acurácia em português é de 85,0\% em A e de 82,7\% em B.

Quanto ao aspecto, a atribuição não depende do modelo de sentimento. Acerta 54,5\% na amostra (kappa 0,50). Ponderando apenas pela participação de cada aspecto no corpus, a estimativa é de 52\%. Com os estratos de aspecto e polaridade, é de 66\% (IC 95\% por \textit{bootstrap} de 57\% a 75\%), uma estimativa instável por causa de células com poucas sentenças e pesos muito desiguais. Na população, a acurácia de aspecto situa-se, portanto, entre 52\% e 66\%. O acerto é parecido nas três polaridades previstas (54\%, 53\% e 57\%), porém muito diferente conforme o tópico: 65\% nas sentenças de tópico direto (186) e 38\% nas de tópico reatribuído por similaridade (111). A Tabela~\ref{tab:precisao} mostra a precisão por aspecto, com cerca de 27 sentenças rotuladas em cada um (intervalos de aproximadamente $\pm 0{,}18$).

\begin{table}[htb]
\centering
\caption{Precisão da atribuição de aspecto, por aspecto previsto (27 sentenças rotuladas em cada).}
\small
\begin{tabular}{l|r}
\hline
Aspecto previsto & Precisão\\
\hline
Custo-benefício & 74\%\\
Localização & 74\%\\
Precisão do anúncio & 70\%\\
Limpeza & 67\%\\
Check-in & 56\%\\
Segurança & 48\%\\
Comodidades & 44\%\\
Comunicação & 44\%\\
Ruído e conforto & 44\%\\
Satisfação geral & 41\%\\
Fora do perfil (\texttt{geral}) & 37\%\\
\hline
\end{tabular}
\label{tab:precisao}
\end{table}

\subsection{Concordância entre os dois modelos}

Com a evidência direta limitada a 297 sentenças, o corpus inteiro permite comparar os modelos entre si. Nas 1.512.063 sentenças, a correlação de postos das estrelas esperadas de A e B é 0,61 (Pearson 0,66) e o kappa ponderado linear das polaridades é 0,45. Por idioma, o kappa é 0,40 em português, 0,54 em inglês e 0,54 em espanhol. No nível de anúncio e aspecto, a concordância é maior (correlação de postos de 0,69 e kappa de 0,55), porque a média de várias sentenças reduz o ruído. Ela varia de 0,46 (precisão do anúncio) a 0,86 (comodidades). Os dois modelos ordenam sentenças e anúncios de modo semelhante, mas não são intercambiáveis: das 1.192.861 sentenças que A classifica como positivas, B concorda em 1.087.384. Das 215.635 que A classifica como neutras, B chama 128.723 de positivas.

\subsection{Regra de escolha fixada de antemão e modelo principal}

A regra de escolha foi registrada antes da execução do modelo B (Seção 3). A diferença B menos A no kappa ponderado linear populacional é de $-0{,}141$ (IC 95\% de $-0{,}272$ a $-0{,}012$). O intervalo exclui zero, mas a favor de A e por isso o modelo principal é A. Na acurácia populacional a diferença é de $-0{,}023$ (IC de $-0{,}078$ a $0{,}021$). Apenas em português, no kappa ponderado linear, a diferença é de $-0{,}169$ (IC de $-0{,}323$ a $0{,}011$), intervalo que inclui zero. A escolha não é óbvia, porque B é melhor na amostra e A na população. A regra foi fixada de antemão justamente para não escolher o resultado mais conveniente. Daqui em diante o modelo principal é A. Os resultados de B acompanham cada análise como verificação de robustez.

\subsection{Validade de construto: concordância com as subnotas oficiais}

A evidência indireta mais importante é a concordância com as subnotas dadas pelos próprios hóspedes. Para cada aspecto com subnota correspondente, correlacionou-se, no nível do anúncio, o sentimento do aspecto com a subnota e a satisfação geral com a nota global. Não há limiar de aprovação e interessam a força e o sinal da correlação. A Tabela~\ref{tab:construto} apresenta os resultados.

\begin{table}[htb]
\centering
\caption{Correlação entre o sentimento do aspecto e a subnota oficial correspondente, no nível do anúncio.}
\small
\begin{tabular}{l|l|r|r|r|r|r}
\hline
Aspecto & Subnota & $n$ & Spearman (A) & Pearson (A) & Spearman (B) & Pearson (B)\\
\hline
Limpeza & \texttt{cleanliness} & 19.095 & 0,45 & 0,58 & 0,42 & 0,52\\
Localização & \texttt{location} & 24.101 & 0,24 & 0,31 & 0,24 & 0,27\\
Check-in & \texttt{checkin} & 8.919 & 0,42 & 0,40 & 0,43 & 0,42\\
Comunicação & \texttt{communication} & 25.418 & 0,30 & 0,44 & 0,30 & 0,43\\
Custo-benefício & \texttt{value} & 3.552 & 0,24 & 0,26 & 0,22 & 0,23\\
Precisão do anúncio & \texttt{accuracy} & 16.638 & 0,17 & 0,22 & 0,15 & 0,21\\
Satisfação geral & nota global & 26.378 & 0,34 & 0,52 & 0,26 & 0,38\\
\hline
\end{tabular}
\label{tab:construto}
\end{table}

Todos os coeficientes são positivos e significativos nos dois modelos. A concordância é moderada, de 0,17 a 0,45 em A e de 0,15 a 0,43 em B, sem um modelo sistematicamente melhor: B é um pouco maior em check-in e comunicação; A é maior em limpeza, custo-benefício e, sobretudo, satisfação geral. As correlações devem ser lidas à luz de que as subnotas também são comprimidas, com muitos empates em 5,0, o que as atenua. Cinco verificações de robustez complementam o resultado.

Na primeira, por volume de avaliações, a concordância dos anúncios com até 30 avaliações (em que o sentimento usa todas as avaliações) difere da dos anúncios com mais de 30 em no máximo 0,11 (check-in), sem padrão sistemático. A segunda trata da validade convergente e discriminante: a subnota correspondente é a mais correlacionada com o aspecto em 5 de 7 casos em A e em 3 de 7 em B. Nos demais casos fica em segundo ou terceiro lugar, porque as subnotas são muito correlacionadas entre si (efeito halo). Controlando a nota global por correlação parcial, resta informação específica em limpeza (0,17), localização (0,22), check-in (0,23) e comunicação (0,15). Em custo-benefício (0,08) e precisão do anúncio (0,04) resta quase nenhuma. Na terceira, a probabilidade de o sentimento ser maior nos anúncios de subnota 5 do que nos de subnota menor (AUC) vai de 0,57 a 0,69 em A e de 0,55 a 0,67 em B. Na quarta, as correlações são positivas em português, inglês e espanhol (de 0,13 a 0,45 em A) e em casa ou apartamento inteiro e quarto privado (de 0,17 a 0,46), o que sustenta o uso de modelos multilíngues. Na quinta, a confiabilidade das metades do sentimento por aspecto, com correção de Spearman-Brown, é baixa, de 0,10 a 0,40 em A e de 0,06 a 0,40 em B. As correlações corrigidas por essa atenuação (de 0,57 a 0,96 em A) são apenas ilustrativas, porque com confiabilidades tão baixas a correção é instável e corrige só o lado do sentimento.

\subsection{Checagem com um segundo método de tópicos e com um segundo anotador}

Para saber se os resultados por aspecto dependem do método de tópicos, cada sentença foi atribuída ao tópico dominante de um modelo LDA com 10 tópicos. O mapa dos tópicos do LDA para aspectos seguiu uma regra: o tópico recebe aspecto apenas quando a leitura das palavras de maior peso, o aspecto mais similar aos termos iniciais e uma margem de 0,05 entre o primeiro e o segundo aspecto coincidem; caso contrário, o tópico é considerado misto e fica fora do perfil. O mapa foi confirmado pelo pesquisador. Quatro dos dez tópicos receberam aspecto (comodidades, comunicação, satisfação geral e localização). Além disso, 55\% das sentenças caem em tópicos mistos, de modo que o LDA permite checar apenas esses quatro aspectos. Entre as sentenças com aspecto nos dois métodos, 68,8\% recebem o mesmo aspecto. A Tabela~\ref{tab:lda} compara os resultados.

\begin{table}[htb]
\centering
\caption{Checagem com LDA: correlação do perfil, validade de construto e parcela do $R^2$ da nota global.}
\small
\begin{tabular}{l|r|r|r|r|r}
\hline
Aspecto & Perfil BERTopic $\times$ LDA & Subnota (LDA) & Subnota (BERTopic) & $R^2$ (LDA) & $R^2$ (BERTopic)\\
\hline
Localização & 0,61 & 0,20 & 0,24 & 1,6\% & 0,9\%\\
Comunicação & 0,54 & 0,28 & 0,30 & 22,7\% & 15,5\%\\
Comodidades & 0,67 & -- & -- & 62,3\% & 59,0\%\\
Satisfação geral & 0,46 & 0,19 & 0,34 & 13,4\% & 24,6\%\\
\hline
\end{tabular}
\label{tab:lda}
\end{table}

O sinal da validade de construto se mantém nos três aspectos que têm subnota, com correlações do mesmo tamanho nos dois mais específicos e menores em satisfação geral. Na decomposição do $R^2$ da nota global, comodidades fica em primeiro nos dois métodos; satisfação geral e comunicação trocam de posição (correlação de postos de 0,80, com apenas quatro aspectos).

Um segundo anotador rotulou, sem ver as classificações do primeiro nem as dos modelos, 100 sentenças da amostra. A concordância no sentimento é de 91\% (kappa de Cohen de 0,86, IC de 0,76 a 0,94; kappa ponderado linear de 0,88) e, no aspecto, de 67\% (kappa de 0,62, IC de 0,51 a 0,73). Os anotadores concordam bem quanto ao sentimento e de forma moderada quanto ao aspecto. A concordância é baixa justamente em satisfação geral (19\%, $n = 16$, em geral confundida com o grupo \texttt{geral}) e em limpeza (33\%, $n = 6$) e alta em custo-benefício e precisão do anúncio (100\% nos dois). A amostra é pequena e as concordâncias por aspecto se apoiam em 4 a 20 sentenças.

\subsection{Classificação dos aspectos}

A partir das evidências, cada aspecto foi classificado por um critério descritivo, fixado antes das seções seguintes. O aspecto é considerado confiável quando a precisão de atribuição na amostra é de ao menos 0,60 e a correlação com a subnota correspondente é de ao menos 0,20 nos dois modelos. Nos demais casos pede cautela. A Tabela~\ref{tab:classif} apresenta o resultado.

\begin{table}[htb]
\centering
\caption{Classificação dos aspectos quanto à confiança da medida.}
\small
\begin{tabular}{l|r|r|r|l}
\hline
Aspecto & Precisão & Spearman (A) & Spearman (B) & Classificação\\
\hline
Limpeza & 67\% & 0,45 & 0,42 & confiável\\
Localização & 74\% & 0,24 & 0,24 & confiável\\
Custo-benefício & 74\% & 0,24 & 0,22 & confiável\\
Check-in & 56\% & 0,42 & 0,43 & cautela\\
Comunicação & 44\% & 0,30 & 0,30 & cautela\\
Precisão do anúncio & 70\% & 0,17 & 0,15 & cautela\\
Satisfação geral & 41\% & 0,34 & 0,26 & cautela\\
Ruído e conforto & 44\% & -- & -- & cautela (sem subnota)\\
Segurança & 48\% & -- & -- & cautela (sem subnota)\\
Comodidades & 44\% & -- & -- & cautela (sem subnota)\\
\hline
\end{tabular}
\label{tab:classif}
\end{table}

São confiáveis limpeza, localização e custo-benefício. Os sete restantes pedem cautela. Duas consequências valem para as seções seguintes. O aspecto de maior peso na Seção 5, comodidades, está entre os que pedem cautela, por ter atribuição imprecisa e nenhuma subnota que o valide. Por isso a leitura prática se apoia nos tópicos finos de queixa e elogio, cuja atribuição é melhor (65\% de acerto de aspecto nos tópicos diretos). Os pesos por aspecto são apresentados como indicativos. Além disso, nenhum resultado por aspecto deve ser lido sem a classificação correspondente.
